\section{Read Completion Tracking and Serialization}
\label{sec:read-completion}

Read completion tracking, final-beat indication, and release are implemented
by the read ROB instance and its Retiring Register (\S\ref{sec:rob}).

\subsection{Read ROB metadata and completion}

The read Metadata RAM is directly indexed by \texttt{ROB\_INDEX}.  Each entry
contains \texttt{exp\_pkts[6:0]} (the number of distinct 64-B lines the read
burst touches, \texttt{1..64}), a 64-bit completion bitmap, and
\texttt{rd\_dbuf\_base}, the base of a contiguously-allocated RD\_SRAM range
of \texttt{exp\_pkts} 64-B lines.  \texttt{exp\_pkts} is initialized from the
Burst Splitter --- one MC packet per distinct 64-B line, not per AXI beat
(\S\ref{sec:splitter}) --- and is the exact MC read-line count.  The read
entry stores no AXI address or 64\,B line footprint for hazard purposes:
same-line read/write ordering is enforced entirely at read admission, before
allocation (\S\ref{sec:samelinehaz}), and a completed read is never
re-checked.

\textbf{Read Serializer dependency.}  Reconstructing the AXI R-channel beat
stream (\S\ref{sec:read-serializer}) needs retained AXI burst metadata ---
\texttt{ARLEN}, \texttt{ARSIZE}, \texttt{ARBURST}, and a normalized base /
\texttt{WrapBoundary} --- in addition to \texttt{exp\_pkts} and
\texttt{rd\_dbuf\_base}.  This metadata is a required input to retirement;
its exact storage location (Metadata RAM, a side structure, Slot-FIFO
context, or other) is \textbf{subject to ROB reconciliation} and is not
fixed here.

On a read completion, \texttt{GLOBAL\_TAG[13:0]} directly selects the entry
(\texttt{[13:6]}) and its packet-number bit (\texttt{[5:0]}).  The returned
64-B data unit is written to \texttt{rd\_dbuf\_base + PACKET\_NUMBER} in
RD\_SRAM.  RD\_SRAM depth (hence \texttt{RD\_SRAM\_ADDR\_W}), its
allocation/read/free handshakes, and latency remain \textbf{TBD}; the ROB
assumes only that the contiguous \texttt{rd\_dbuf\_base} range is valid from
allocation until the contract-defined release point.

\subsection{Read retirement}

An AXID Slot FIFO head can retire only after \texttt{completion\_bitmap}
bits \texttt{[exp\_pkts-1:0]} are all set --- i.e.\ every touched 64-B line
has returned.  The Retiring Register then holds that index while the Read
Serializer (\S\ref{sec:read-serializer}) emits \texttt{ARLEN + 1} AXI R
beats, mapping each beat to the line it is sourced from and reading that
complete line from RD\_SRAM.  \texttt{rob\_last} is generated with the final
AXI beat (\texttt{beat\_index == ARLEN}), not the final line.  Final AXI R
acceptance, after the \texttt{rd\_dbuf\_base} range is released, permits
retirement to pop the Slot FIFO and return \texttt{ROB\_INDEX} to the read
Free List.  Completion alone never frees the index.

\subsection{Read Serializer}
\label{sec:read-serializer}

The Read Serializer reconstructs the original AXI R-channel beat stream from
the returned 64-B lines while the Retiring Register holds the read
transaction.  The behavior below is architectural; the storage of its AXI
metadata inputs is subject to ROB reconciliation (\S\ref{sec:read-completion}).

\begin{itemize}[noitemsep]
  \item Maintains an AXI beat counter \texttt{beat\_index = 0 .. ARLEN} and
        emits \texttt{ARLEN + 1} R beats in order.
  \item For each beat it forms the AXI beat byte address per
        \texttt{ARBURST}/\texttt{ARSIZE}: \texttt{INCR} increments by
        \texttt{NB = 1 << ARSIZE} with the AXI4 unaligned-first-transfer
        rule; \texttt{WRAP} increments modulo \texttt{Total} relative to
        \texttt{WrapBoundary}; \texttt{FIXED} repeats \texttt{ARADDR}.
  \item Maps the beat address to a line index
        \texttt{PACKET\_NUMBER = (beat\_addr >> 6) - first\_line}
        (\texttt{0..exp\_pkts-1}) and reads the complete 64-B line from
        \texttt{rd\_dbuf\_base + PACKET\_NUMBER} in RD\_SRAM.
  \item Drives \texttt{RDATA} directly from that line.  AXI byte lanes and
        64-B line offsets are congruent, so no barrel shifter or data
        realignment is used; the addressed master selects its byte lanes.
  \item Consecutive beats may map to the same line index and re-use the same
        RD\_SRAM line.
  \item \texttt{RRESP} for a beat is the per-line status of the line it is
        sourced from (\S\ref{sec:cifmc}); beats sourced from the same line
        carry the same \texttt{RRESP}.
  \item \texttt{rob\_last}/\texttt{RLAST} is asserted on
        \texttt{beat\_index == ARLEN} and accepted downstream before the
        \texttt{rd\_dbuf\_base} range is released and the entry retires.
\end{itemize}

The line-set parameters \texttt{first\_line}, \texttt{NB}, \texttt{Total},
and \texttt{WrapBoundary} are the ones computed by Burst Splitter Stage~1
(\S\ref{sec:splitter}).
